% The error budget, drawn as what it is: a stack of terms, most of which this
% bench cannot reduce. Term names sit outside the bars, because a bar whose
% width is the quantity cannot also be a label.
\begin{tikzpicture}[font=\sffamily\scriptsize, x=1cm, y=1cm]

\tikzset{
  errbar/.style={draw=linecol, anchor=south west, minimum height=5mm, inner sep=0pt},
  fixable/.style={errbar, fill=usrcol},
  stuck/.style={errbar, fill=hwcol},
  unknown/.style={errbar, fill=black!10},
  term/.style={font=\sffamily\scriptsize, anchor=east, align=right, text width=40mm},
  why/.style={font=\sffamily\tiny, anchor=west, text=black!60, align=left},
}

\node[chlabel, anchor=west] at (-4.3,4.5)
  {every term in the estimate's error, and whether this bench can do anything about it};

\foreach \i/\w/\sty/\lab/\note in {%
  0/3.4/fixable/{position quantisation, doubly differentiated}/{a finer encoder, or a longer filter},
  1/2.6/stuck/{filter latency}/{trades directly against the term above},
  2/3.0/stuck/{inertia parameter error}/{J was chosen; there is nothing here to weigh},
  3/4.2/stuck/{friction model error}/{largest at low speed and at reversals},
  4/3.8/stuck/{commanded torque error}/{no current loop: an assumption twice over},
  5/1.2/stuck/{one control period of latency}/{fixed by the period},
  6/4.6/unknown/{unmodelled dynamics}/{backlash, compliance, cogging: size unknown}}{
  \node[term] at (-0.2,{4.0-\i*0.62}) {\lab};
  \node[\sty, minimum width={\w*10mm}] at (0,{3.78-\i*0.62}) {};
  \node[why] at ({\w+0.15},{4.0-\i*0.62}) {\note};
}

% ---------------- the legend
\node[fixable, minimum width=10mm] at (0,-0.75) {};
\node[font=\sffamily\tiny, anchor=west] at (1.15,-0.5) {this bench can reduce it};
\node[stuck, minimum width=10mm] at (4.0,-0.75) {};
\node[font=\sffamily\tiny, anchor=west] at (5.15,-0.5) {it cannot};
\node[unknown, minimum width=10mm] at (7.0,-0.75) {};
\node[font=\sffamily\tiny, anchor=west] at (8.15,-0.5) {size unknown};

\node[chlabel, anchor=north west, text width=52mm] at (-4.3,-1.3)
  {Bar widths are indicative of relative size rather than measured, and the
   figure says so: a measured error budget for this estimate would require the
   sensor whose absence is the subject of the chapter.};

\node[umlnote, text width=56mm, anchor=north west] at (2.6,-1.3)
  {Four terms this bench cannot reduce and one whose size is unknown. That is
   the honest summary of the estimate, and it is why the detector built on it is
   tuned by its measured false-alarm rate rather than by a threshold in newton
   metres that would imply a precision the number does not have.};

\node[chlabel, anchor=north west, text width=116mm] at (-4.3,-4.2)
  {The arithmetic behind the first bar, which is worth carrying: with four
   thousand counts per revolution and a one millisecond period, one count of
   position quantisation becomes roughly 1.6 radians per second squared of
   apparent acceleration, and the modelled inertia turns that straight into an
   apparent torque. That is the noise floor before any filtering, and the filter
   that reduces it is the same filter that delays the detection.};
\end{tikzpicture}
